\documentclass[11pt,a4paper]{article}
\usepackage[utf8]{inputenc}
\usepackage[T1]{fontenc}
\usepackage{geometry}
\geometry{left=2.7cm,right=2.7cm,top=2.6cm,bottom=2.6cm}
\usepackage{lmodern,microtype,setspace}
\usepackage{graphicx,float,amsmath,amssymb,cite,listings,xcolor,booktabs,array,caption,subcaption,algorithm,algpseudocode,tikz,pgfplots,seqsplit}
\usepackage[hidelinks]{hyperref}
\pgfplotsset{compat=1.18}
\usetikzlibrary{shapes,arrows,positioning,fit,backgrounds,matrix,calc}
\onehalfspacing
\setlength{\parindent}{1.5em}
\setlength{\parskip}{0.15em}
\hypersetup{
  pdftitle={Privacy Lens: An Evidence-Bounded Android Privacy Review Framework, Revision 23},
  pdfauthor={Zhiheng Zhang}
}
\title{Privacy Lens: An Evidence-Bounded Android Privacy Review Framework with Offline-Verified Legal Source Content (Revision 23)}
\author{Zhiheng Zhang \\ Student ID: 54560484 \\ Supervisor: Richard Sinnott}
\date{August 2026}
\begin{document}
\maketitle
\begin{abstract}
\small\singlespacing
Mobile privacy-review tools face an evidence-to-decision gap: a permission observation can indicate activity but cannot by itself establish purpose, necessity, lawful basis, acquisition completeness, or legal compliance. This thesis addresses that gap through Privacy Lens, a local-first Android prototype that converts bounded permission-audit summaries into traceable prompts for human review.

Following a design-science method, the architecture treats provenance, temporal scope, missing context, and uncertainty as first-class evidence rather than converting a count into a legal verdict. A fail-closed engine rejects malformed records and unsafe temporal aggregation, isolates simulator data, and records the rule-pack version applied to each finding. Jurisdiction mappings remain outside the Android, repository, and interface layers through a replaceable pack contract; the European Union General Data Protection Regulation provides the evaluated legal objective.

Revision 23 closes one additional provenance gap by defining a canonical source-content manifest whose entries bind each recorded official source to an exact URL, byte length, retrieval date, and SHA-256 digest. A reusable offline verifier recomputes these values from the six retrieved official documents and the runtime gate requires every expected artifact byte sequence to verify before reassurance. The legal pack deliberately remains unattested, its source bytes are not bundled with the application, and the production trust store remains empty. Findings remain local; backup and remote updates are disabled; the evaluated package has no network or sensitive runtime permission; and users can clear evidence directly. The interface separates signal, missing context, proportionate action, source-content verification, review signature, and production key governance.

In a fixed-seed 1,000-case corpus, the engine matched its specified oracle with 800 true positives and 200 true negatives. A further 1,800 independently calculated pack cases, manifest coverage and ordering probes, missing, length-mismatched, and same-length hash-mismatched artifact probes, signature and key-governance probes, accessibility contracts, and release-privacy contracts passed. All six retrieved official documents verified against the manifest and a one-byte mutation failed closed. The 1.11.0 APK and AAB built successfully. On one OPPO PERM00, the interface exposed the three-link evidence chain and the absent source-byte state; five final-artifact cold starts returned status ok with an empty crash buffer.

These results support an initial store-submission engineering candidate within the recorded scope. A matching SHA-256 digest supports byte identity relative to the recorded retrieval but does not prove publisher identity, legal status, completeness, or correctness. Signature verification authenticates only a payload relative to a configured key. The evidence does not establish reviewer identity or qualification, legal correctness, unrestricted Android visibility, accessibility conformance, participant outcomes, long-run reliability, store approval, or publication acceptance. The contribution is an evidence-bounded architecture and traceability method that separate statutory motivation, engineering heuristics, observed evidence, missing legal facts, and prohibited conclusions.
\end{abstract}
\newpage
\tableofcontents
\newpage
\input{Iteration-V2/Chapter1/chapter-01-23}
\input{Iteration-V2/Chapter2/chapter-02-23}
\input{Iteration-V2/Chapter3/chapter-03-23}
\input{Iteration-V2/Chapter4/chapter-04-23}
\input{Iteration-V2/Chapter5/chapter-05-23}
\input{Iteration-V2/Chapter6/chapter-06-23}
\input{Iteration-V2/Chapter7/chapter-07-23}
{\small
\bibliographystyle{ieeetr}
\bibliography{reference}
}
\end{document}
